\chapter{Les données}


\section{Constitution de la base de données}
\label{sec:data}

\subsection{Source}

Les données ont été extraites du site \texttt{football-data.co.uk}, couvrant
dix saisons consécutives (2014-2024) des cinq grands championnats européens :
Premier League (Angleterre), LaLiga (Espagne), SerieA (Italie), Bundesliga
(Allemagne) et Ligue1 (France). Cela représente donc 17\,786 matchs.
Après nettoyage, la répartition des résultats est la suivante :
\begin{itemize}
    \item Victoires à domicile : 7\,905 matchs (45\,\%) ;
    \item Matchs nuls : 5\,420 matchs (25\,\%) ;
    \item Victoires à l'extérieur : 4\,461 matchs (30\,\%).
\end{itemize}

Cette distribution confirme l'avantage du terrain dans le football abordé par Dixon \& Coles.

\subsection{Structure des données brutes}

Les colonnes des données brutes sont organisées comme suit.

\paragraph{Informations générales.}
\begin{center}
    \begin{tabular}{ll}
        \toprule
        \textbf{Colonne}  & \textbf{Description}       \\
        \midrule
        \texttt{Div}      & Division (championnat)     \\
        \texttt{Season}   & Saison                     \\
        \texttt{League}   & Nom de la ligue            \\
        \texttt{Date}     & Date du match (aaaa/mm/jj) \\
        \texttt{HomeTeam} & Équipe à domicile          \\
        \texttt{AwayTeam} & Équipe à l'extérieur       \\
        \bottomrule
    \end{tabular}
\end{center}

\paragraph{Résultats et statistiques de match.}
Les variables de résultat et de jeu couvrent la fin de match
(\textit{Full Time}) et la mi-temps (\textit{Half Time}) :
\begin{itemize}
    \item \texttt{FTHG} / \texttt{FTAG} : buts marqués (domicile / extérieur) ;
    \item \texttt{FTR} : résultat final (1 : domicile, 2 : extérieur, 0 : nul) ;
    \item \texttt{HTHG} / \texttt{HTAG} / \texttt{HTR} : équivalents à la mi-temps ;
    \item \texttt{HS} / \texttt{AS} : tirs ; \texttt{HST} / \texttt{AST} : tirs cadrés ; \texttt{HC} / \texttt{AC} : corners ;
    \item \texttt{HF} / \texttt{AF} : fautes ; \texttt{HY} / \texttt{AY} : cartons jaunes ; \texttt{HR} / \texttt{AR} : cartons rouges.
\end{itemize}

% ============================================================


\noindent L'étape de préparation a été réalisée sous Python avec la bibliothèque
\texttt{Pandas}. Elle consiste à transformer les données brutes issues des
dix saisons et cinq championnats en un dataset exploitable pour l'apprentissage.

% ============================================================


\section{Ingénierie des descripteurs}
\label{sec:features}

Pour enrichir la capacité prédictive, 61 caractéristiques ont été construites,
réparties en plusieurs groupes fonctionnels.

\subsection{Indicateurs de forme et de dynamique}

\paragraph{Fenêtre glissante.}
La moyenne des buts (forme), des tirs tentés (attaque) et des buts subis (défense)
est calculée sur les 5 derniers matchs de chaque équipe, aussi bien pour les
statistiques de fin de match (\textit{Full Time}) que de mi-temps (\textit{Half Time}),
afin de capturer l'état de forme récent.

\paragraph{Précision et efficacité offensive.}
Le rapport $\frac{\text{tirs cadrés}}{\text{tirs tentés}}$ mesure la qualité offensive. Une variante
pondérée (\texttt{FT\_prec\_weight}) applique un facteur $w_g = 2$ aux tirs
convertis en buts (nous considérons qu'un but est deux fois plus important qu'un 
tir cadré).

\paragraph{Discipline.}
Les moyennes de fautes et de cartons (jaunes et rouges) sur les 5 derniers matchs
permettent de capturer l'agressivité et le risque d'infériorité numérique.

\paragraph{Séries consécutives (\textit{Streaks}).}
Deux compteurs de série sont maintenus pour chaque équipe :
\begin{itemize}
    \item \texttt{H\_WinStreak} / \texttt{A\_WinStreak} : nombre de victoires consécutives avant le match ;
    \item \texttt{H\_LoseStreak} / \texttt{A\_LoseStreak} : nombre de défaites consécutives correspondantes.
\end{itemize}
La mise à jour respecte l'absence de \textit{data leakage} : la valeur enregistrée
pour un match est celle du compteur \textit{avant} la rencontre. Tout match nul
remet les deux compteurs à zéro.

\paragraph{Repos.}
Les variables \texttt{HRepos} et \texttt{ARepos} mesurent le nombre de jours depuis
le dernier match (plafonné à 25 jours), capturant la fatigue potentielle ou la fraîcheur physique.

\subsection{Modélisation du niveau relatif}

\paragraph{Score Elo.}
Le score Elo est mis à jour itérativement après chaque match avec un facteur de
pondération $G$ lié à l'écart de buts, constituant la base de l'évaluation du niveau des équipes.

\paragraph{Classement ('Rank').}
Nous avons également calculé un classement basé sur les points au cours d'une saison.
Une victoire rapporte 3 points, un match nul 1 point, et une défaite 0 point.
Ce classement est mis à jour après chaque match, reflétant la performance relative des équipes au fil de la saison.

\paragraph{Variables de comparaison.}
En s'inspirant des travaux de Dixon \& Coles\cite{dixon} et Egidi \&
Torelli\cite{egidi}, des variables de comparaison entre l'équipe à domicile
(H) et l'équipe visiteuse (A) sont générées pour chaque groupe de descripteurs selon deux formules :

\begin{center}
    \begin{tabular}{lll}
        \toprule
        \textbf{Suffixe} & \textbf{Calcul} & \textbf{Justification théorique}           \\
        \midrule
        \texttt{\_diff}  & $H - A$         & Écart de niveau explicite \cite{egidi} \\
        \texttt{\_ratio} & $H / A$         & Structure multiplicative \cite{dixon}  \\
        \bottomrule
    \end{tabular}
\end{center}

La formule du ratio utilise $\frac{stat\_H + 1}{stat\_A +1}$ pour éviter les 
divisions par zéro.

\subsection{Récapitulatif des indicateurs calculés}

\begin{longtable}{lllp{5.5cm}}
    \toprule
    \textbf{Catégorie} & \textbf{Domicile} & \textbf{Extérieur} & \textbf{Description} \\
    \midrule
    \endfirsthead
    \toprule
    \textbf{Catégorie} & \textbf{Domicile} & \textbf{Extérieur} & \textbf{Description} \\
    \midrule
    \endhead
    Forme (FT/HT)     & \texttt{XX\_Hforme}       & \texttt{XX\_Aforme}       & Moyenne buts (sur 5 matchs) \\
    Attaque (FT)      & \texttt{FT\_Hatt}         & \texttt{FT\_Aatt}         & Moyenne des tirs effectués \\
    Défense (FT/HT)   & \texttt{XX\_Hdef}         & \texttt{XX\_Adef}         & Moyenne des tirs concédés \\
    Elo               & \texttt{XX\_Elo\_H}        & \texttt{XX\_Elo\_A}        & Classement elo avant match \\
    Rank              & \texttt{H\_Rank}         & \texttt{A\_Rank}         & Classement basé sur les points  \\
    Précision         & \texttt{FT\_Hprecision}   & \texttt{FT\_Aprecision}   & Ratio tirs cadrés / tirs totaux \\
    Efficacité pond.  & \texttt{FT\_Hprec\_weight} & \texttt{FT\_Aprec\_weight} & Précision pondérée ($w_g = 2$) \\
    Fautes            & \texttt{FT\_HavgF}        & \texttt{FT\_AavgF}        & Moyenne fautes (5 matchs) \\
    Cartons J.        & \texttt{FT\_HavgY}        & \texttt{FT\_AavgY}        & Moyenne jaunes (5 matchs) \\
    Cartons R.        & \texttt{FT\_HavgR}        & \texttt{FT\_AavgR}        & Moyenne rouges (5 matchs) \\
    Repos             & \texttt{HRepos}          & \texttt{ARepos}          & Jours depuis dernier match (max 25) \\
    \bottomrule
    \caption{Indicateurs de performance calculés par équipe. (XX=FT ou HT)}
\end{longtable}

% ============================================================
\section{Transformation et Normalisation des Données des bookmakers}
Afin d'établir une comparaison entre notre modèle de classification 
binaire et les bookmakers, nous avons procédé à une restructuration des 
cotes des bookmakers (Bet365, Pinnacle, etc.). La méthodologie employée est la suivante, 
dans un premier temps inversion des cotes afin de convertir les cotes décimales en 
probabilités implicites brutes selon la formule $P_{brute} = \frac{1}{\text{cote}}$.
Par la suite, afin de pouvoir comparer avec un modèle OvR, nous avons binarisé les issues 
du match nul ($P_D$) et de la victoire extérieure ($P_A$) pour créer la classe 
complémentaire au succès à domicile, notée $P_{Reste}$. Enfin, les probabilités brutes 
ont été normalisées afin que la somme des probabilités de chaque 
bookmaker soit égale à 1.

Cette étape garantit que le modèle et les bookmakers sont évalués sur le même 
espace probabiliste : $P(Home) + P(Reste) = 1$.


Par la suite, nous avons procédé à un classement des bookmakers 
entre eux à l'aide d'un Critical Difference diagramme\footnote{https://mirkobunse.github.io/CriticalDifferenceDiagrams.jl/dev/}. 
Cette méthode se 
base sur un test de Friedman afin de déterminer s'il existe ou non des différences 
significative entre les éléments comparer. Une fois ce test réalisé, l'analyse peut 
avoir lieu. 

\begin{figure}[H]
    \centering 
    \includegraphics[width=0.75\textwidth]{../../plot/cdd.png}
    \caption{Diagram de différences critique entre les bookmaker}
    \label{fig:cdd}
\end{figure}
\noindent Sur la figure \ref{fig:cdd} le trait noir reliant deux couleur indique qu'il n'y a pas 
de différence significative entre les deux bookmakers. 



\section{Définition des variables cibles}
\label{sec:labels}

Le problème est traité comme une classification multi-classes.
Pour faciliter l'apprentissage et permettre une approche One-vs-Rest,
trois colonnes cibles binaires sont créées :

\begin{itemize}
    \item \texttt{Hvs} : victoire à domicile vs (nul ou défaite) ;
    \item \texttt{Avs} : victoire à l'extérieur vs (nul ou victoire à domicile) ;
    \item \texttt{Dvs} : match nul vs (victoire ou défaite).
\end{itemize}

Cela permet de rééquilibrer les classes, nous passons d'une distribution 
45\% - 25\% - 30\% à 45\% - 55\%.
De plus, par la suite, si nous souhaitons entrainer un modèle par cible,
cela nous permettra si nécessaire de combiner les modèles afin
d'obtenir une prédiction plus robuste.


\subsection{Validation Statistique par Tests de Friedman et Nemenyi}

L'évaluation de la précision des bookmakers repose sur une approche
statistique rigoureuse en deux étapes, visant à comparer les rangs
moyens des \textit{Log-Loss} binaires obtenues sur l'issue « Home
vs Reste ».

\begin{itemize}
    \item \textbf{Test de Friedman :} Nous utilisons ce test
 non-paramétrique pour vérifier s'il existe une différence
 significative globale entre les performances des différents
 bookmakers. Le rejet de l'hypothèse nulle ($p < 0,05$) indique
 qu'au moins un bookmaker se distingue statistiquement des autres.
    \item \textbf{Test Post-hoc de Nemenyi :} Suite au test de Friedman,
 le test de Nemenyi est appliqué pour identifier les paires
 de bookmaker dont les performances diffèrent significativement.
 Ce test permet de calculer la \textit{Critical Difference} (CD).
    \item \textbf{Interprétation du CD Diagram :} Sur le diagramme généré,
 la position sur l'axe des abscisses indique le rang moyen (un
 rang proche de 1 étant optimal). Les segments horizontaux relient
 les groupes de bookmakers dont la différence de performance n'est pas
 statistiquement significative au seuil $\alpha = 0,05$.
\end{itemize}



% \subsubsection{Détails mathématiques et interprétation des p-values}

% L'analyse s'appuie sur une hiérarchie de tests statistiques garantissant
% la robustesse des conclusions :

% \begin{enumerate}
%     \item \textbf{Test de Friedman :} Ce test transforme les scores de
%      \textit{Log-Loss} en rangs pour chaque observation. La statistique
%  calculée suit approximativement une loi du $\chi^2$. Une \textit
% {p-value} globale inférieure à 0,05 nous permet de rejeter l'hypothèse
%  nulle d'équivalence des performances. Ce test est privilégié ici
%  car il ne nécessite pas l'hypothèse de normalité des distributions
%  de pertes, souvent violée dans les données de paris sportifs.

%     \item \textbf{Test Post-hoc de Nemenyi :} Ce test effectue des
%  comparaisons par paires. Il définit une distance seuil appelée
%  \textit{Critical Difference} (CD). Pour chaque paire de bookmakers,
%  une \textit{p-value} locale est calculée :
%     \begin{itemize}
%         \item Si $p < 0,05$, l'écart entre les rangs moyens est supérieur à
%      la CD, marquant une supériorité statistique.
%         \item Si $p \geq 0,05$, les acteurs sont regroupés au sein d'une même
%      « clique », matérialisée par un segment horizontal sur le diagramme.
%     \end{itemize}
% \end{enumerate}

% Cette rigueur mathématique permet d'isoler les leaders du marché (ceux situés
% à gauche sans connexions aux acteurs de droite) et de valider si notre modèle
% MLP apporte une valeur ajoutée statistiquement significative par rapport aux
% benchmarks industriels.